% !TEX program = xelatex
% AI-effects 프로젝트 — "한국만을 대상으로 AEI 데이터를 이용해 직업별 시간절감률·
% 직업별 사용량을 계산할 수 있는가"라는 질의에 대해, HuggingFace Anthropic/EconomicIndex
% 저장소의 원자료(특히 release_2026_06_26 V6 CSV 219MB)를 실제로 다운로드해
% geo_id="KOR"/"GLOBAL" 행을 직접 검증한 작업 노트(report/).
% 컴파일: xelatex -> biber -> xelatex -> xelatex (kotex, biblatex/biber 필요)
\documentclass[11pt]{article}

\usepackage{fontspec}
\usepackage{kotex}
\usepackage[a4paper,margin=2.2cm]{geometry}
\usepackage{longtable}
\usepackage{booktabs}
\usepackage{array}
\usepackage{amsmath}
\usepackage{xcolor}
\usepackage{hyperref}
\usepackage{enumitem}
\usepackage{titlesec}
\usepackage[backend=biber,style=authoryear,maxcitenames=2,uniquelist=false,sorting=nyt]{biblatex}
\addbibresource{../../reference/references.bib}

\hypersetup{colorlinks=true,linkcolor=blue!50!black,citecolor=blue!50!black,urlcolor=blue!50!black}
\newcolumntype{L}[1]{>{\raggedright\arraybackslash}p{#1}}
\setlength{\LTleft}{0pt}
\setlength{\LTright}{0pt}
\renewcommand{\arraystretch}{1.15}
\setlength{\tabcolsep}{4pt}

\title{AEI로 한국만을 대상으로\\[1mm]
직업별 시간절감률$\cdot$사용량을 계산할 수 있는가\\[2mm]
\large --- 어떤 릴리스$\cdot$파일을 쓸 것인가, 국가 단위 해상도의 실제 한계, 그리고 GLOBAL 시간절감률을 한국에 적용하는 방식의 타당성 검증 ---}
\author{AI-effects 프로젝트}
\date{\today}

\begin{document}
\maketitle

\begin{abstract}
\noindent 본 노트는 "한국만을 대상으로 Anthropic Economic Index(AEI) 데이터를 이용해 직업별 노동시간 절감률(time savings)과 직업별 AI 사용량을 계산할 수 있는가"라는 질의에 대한 답을 정리한 작업 노트다. \texttt{08\_AEI데이터\_노동시간\_사용량\_스키마.tex}가 AEI 데이터셋의 필드 구조를 일반적으로 정리했다면, 본 노트는 그 구조를 \textbf{한국(\texttt{geo\_id="KOR"})이라는 구체적 국가에 실제로 적용했을 때 무엇이 가능하고 무엇이 불가능한지}를 HuggingFace \texttt{Anthropic/EconomicIndex} 저장소의 원자료를 직접 다운로드해 검증한다. 핵심 결론은 세 가지다. (1) 완성된 형태로 다운로드 가능한 릴리스는 \texttt{release\_2026\_06\_26}(V6) 하나뿐이며, 그 이전 릴리스(V4$\cdot$V5)는 최종 집계본을 공개하지 않는다. (2) 국가 단위(한국 포함)에서는 직업별 시간절감률이 SOC 대분류(Major Group, 22개) 수준까지만 제공되고, SOC 세부직업(Detailed Occupation, 6자리) 수준의 시간절감률은 국가 단위에는 존재하지 않는다---세부직업 수준 시간절감률은 \texttt{GLOBAL} 단위에만 존재한다. (3) 한국의 세부직업별 사용량(424개 O*NET-SOC 코드, 2026년 5월 기준)과 \texttt{GLOBAL} 세부직업별 시간절감률(717개 코드)을 결합하면 커버리지 100\%로 조인되어, 세부직업 단위 분석이 기술적으로 가능하다---단 이는 "한국 고유의 실현 크기"가 아니라 "한국의 사용 비중 $\times$ 전세계(사실상 미국 비중이 큰 표본) 평균 시간절감률"이라는 혼합 지표이므로, 국가 간 시간절감률 동질성이라는 명시적 가정을 전제해야 한다.
\end{abstract}

\tableofcontents

% ============================================================
\section{서론: 질의와 검증 방법}
% ============================================================

\texttt{08\_AEI데이터\_노동시간\_사용량\_스키마.tex}는 AEI 데이터셋의 필드 구조를 릴리스 전체에 걸쳐 문서 수준에서 정리했다. 그러나 "한국만을 대상으로" 분석이 가능한지는 문서만으로는 확정할 수 없다---국가 단위 데이터가 실제로 존재하는지, 프라이버시 임계치를 통과하는지는 원자료를 열어봐야 알 수 있다. 본 노트는 다음 두 질의에 답하기 위해 \texttt{Anthropic/EconomicIndex} 저장소의 \texttt{data\_documentation.md}를 릴리스별로 재대조하고, 실제 CSV 원자료(\texttt{release\_2026\_06\_26/data/aei\_claude\_ai\_2026-06-26.csv}, 219MB)를 다운로드해 \texttt{geo\_id="KOR"}$\cdot$\texttt{geo\_id="GLOBAL"} 행을 직접 집계했다.

\begin{enumerate}[label=(\arabic*)]
  \item 직업별 노동시간 절감률과 직업별 AI 사용량을 한국만을 대상으로 계산하는 것이 가능한가. 가능하다면 어떤 파일$\cdot$릴리스를 써야 하는가.
  \item 직업별 시간절감률은 \texttt{geo\_id="GLOBAL"}의 값을 한국에도 그대로 적용하고, 직업별 사용량은 \texttt{geo\_id="KOR"}의 값을 채택하는 하이브리드 방식이 가능한가. 가능하다면 SOC 분류는 몇 자리까지 가능한가.
\end{enumerate}

% ============================================================
\section{어떤 릴리스$\cdot$파일을 써야 하는가: 완성 집계본이 공개된 릴리스는 V6뿐}
\label{sec:release}
% ============================================================

\texttt{08\_AEI데이터\_노동시간\_사용량\_스키마.tex} 표1은 V4(\texttt{release\_2026\_01\_15})$\cdot$V5(\texttt{release\_2026\_03\_24})도 "동일 스키마 유지"라고만 기술해, 완성된 집계 CSV가 두 릴리스에도 존재하는 것처럼 읽힐 수 있다. 그러나 HuggingFace API(\texttt{tree} 엔드포인트)로 각 릴리스의 \texttt{data/} 폴더를 직접 조회한 결과, 실제 공개 파일 구성은 다음과 같다.

\begin{longtable}{L{2.6cm}L{7.2cm}L{3.6cm}}
\caption{릴리스별 실제 공개 파일 (2026년 9월 HuggingFace API \texttt{tree} 엔드포인트 직접 조회)}\label{tab:files}\\
\toprule
\textbf{릴리스} & \textbf{\texttt{data/} 폴더 실제 내용} & \textbf{완성 집계본 여부}\\
\midrule
\endfirsthead
\multicolumn{3}{l}{\small (표 \thetable\ 계속)}\\
\toprule
\textbf{릴리스} & \textbf{\texttt{data/} 폴더 실제 내용} & \textbf{완성 집계본 여부}\\
\midrule
\endhead
\bottomrule
\endfoot

\texttt{release\_2026\_01\_15} (V4) &
\texttt{data/intermediate/aei\_raw\_claude\_ai\_2025-11-13\_to\_2025-11-20.csv}(94MB, \textbf{1주일치 pre-enrichment 원시본}) &
없음 --- raw 스냅샷만 공개\\

\texttt{release\_2026\_03\_24} (V5) &
\texttt{data/intermediate/aei\_raw\_claude\_ai\_2026-02-05\_to\_2026-02-12.csv}(103MB, \textbf{1주일치 pre-enrichment 원시본}) &
없음 --- raw 스냅샷만 공개\\

\texttt{release\_2026\_06\_26} (V6) &
\texttt{data/aei\_claude\_ai\_2026-06-26.csv}(219MB, \textbf{2026년 4--5월 2개월치, enrichment 완료된 최종본}) &
\textbf{있음}\\

\texttt{labor\_market\_impacts} &
\texttt{job\_exposure.csv}(37KB), \texttt{task\_penetration.csv}(1.9MB) &
있음(단, 지리 컬럼 없음, \S\ref{sec:link} 참조)\\
\end{longtable}

V4$\cdot$V5의 파일명이 \texttt{aei\_raw\_claude\_ai\_*}이고 경로가 \texttt{data/intermediate/}라는 점 자체가 이 파일들이 프라이버시 임계치 적용 이전의 중간 산출물임을 가리킨다. \texttt{data\_documentation.md}(V4)도 "AEI 원시 파일은 원시 건수$\cdot$비율을 담고 있다"고 명시하며, enrichment(지리 코드 3자리 변환, 프라이버시 임계치 적용 등)는 별도 단계로 문서에 기술되어 있을 뿐 그 결과물 파일은 저장소에 공개되어 있지 않다. 반면 V6은 파일명이 \texttt{aei\_claude\_ai\_*}(\texttt{raw} 접두어 없음)이고 \texttt{data/} 최상위에 있어, enrichment가 완료된 배포용 최종본이다.

\textbf{따라서 한국 분석에 쓸 파일은 \texttt{release\_2026\_06\_26/data/aei\_claude\_ai\_2026-06-26.csv} 하나로 좁혀진다.} V4$\cdot$V5의 raw 파일을 쓰려면 연구자가 직접 enrichment 로직(지리 코드 매핑, 최소 관측치 임계치 적용 등)을 재구현해야 하며, 이는 재현성$\cdot$검증 부담이 크므로 권장하지 않는다.

% ============================================================
\section{V6 스키마의 국가 단위 해상도 제약}
\label{sec:schema}
% ============================================================

V6 \texttt{data\_documentation.md}의 컬럼은 \texttt{date\_start, date\_end, geo\_id, geo\_level, category\_name, hierarchy\_level, metric\_id, value, node\_name, node\_external\_id}이며, \texttt{category\_name}은 \texttt{overall}(전체)$\cdot$\texttt{onet}(O*NET 과업)$\cdot$\texttt{request}(요청군집)$\cdot$\texttt{soc\_occupation}(SOC 직업) 4종이다. 이 중 "Metric Availability" 표가 지리 단위별로 어떤 조합에서 전체 지표(시간 지표 포함)가 제공되는지를 명시한다(문서 원문을 그대로 인용).

\begin{longtable}{L{1.8cm}L{2.0cm}L{1.6cm}L{2.6cm}L{2.6cm}L{2.6cm}L{3.4cm}}
\caption{V6 "Metric Availability" 표 (\texttt{release\_2026\_06\_26/data\_documentation.md} 원문)}\label{tab:avail}\\
\toprule
\textbf{Source} & \textbf{Geography} & \textbf{Overall} & \textbf{O*NET} & \textbf{Request} & \textbf{SOC} & \textbf{제공 지표}\\
\midrule
\endfirsthead
\multicolumn{7}{l}{\small (표 \thetable\ 계속)}\\
\toprule
\textbf{Source} & \textbf{Geography} & \textbf{Overall} & \textbf{O*NET} & \textbf{Request} & \textbf{SOC} & \textbf{제공 지표}\\
\midrule
\endhead
\bottomrule
\endfoot

\texttt{claude\_ai} & Global & \checkmark & \checkmark & \checkmark & \checkmark & 전체 지표(시간 지표 포함)\\
\texttt{claude\_ai} & Country & \checkmark & --- & --- & --- & 전체 지표\\
\texttt{claude\_ai} & Country & --- & \checkmark & \checkmark & \checkmark & \texttt{pct}만\\
\texttt{claude\_ai} & Country & --- & GWA(최상위) & Major(최상위) & Major Group(최상위) & 전체 지표\\
\texttt{claude\_ai} & Subregion & \checkmark & --- & --- & --- & 전체 지표\\
\texttt{claude\_ai} & Subregion & --- & \checkmark & \checkmark & \checkmark & \texttt{pct}만\\
\end{longtable}

\noindent 문서 원문: \emph{"A cell is only published if it meets both the aggregation thresholds and the geography sample floor for the given row."}

즉 \textbf{국가(Country) 단위에서는 직업 카테고리(\texttt{soc\_occupation})의 최상위(가장 성긴) 계층---Major Group, SOC 2자리 대분류---에서만 시간 지표를 포함한 전체 지표가 제공}되고, 가장 세밀한 계층(Detailed Occupation, SOC 6자리)에서는 \texttt{pct}(사용 비중)만 제공된다. 세부직업 단위 시간 지표는 오직 \texttt{Global} 행에만 존재한다. 이 표는 \texttt{08\_AEI데이터\_노동시간\_사용량\_스키마.tex}가 정리한 V3$\cdot$V4 시점 문서(국가당 최소 200건 대화, 미국 주당 100건이라는 국가 포함 임계치)와는 별개의, \textbf{셀 단위 해상도 제약}이다.

% ============================================================
\section{한국(\texttt{geo\_id="KOR"}) 실제 데이터 검증}
\label{sec:kor}
% ============================================================

위 문서상 제약이 한국에도 실제로 적용되는지, 그리고 한국이 애초에 국가 포함 임계치를 통과하는지를 219MB 원자료를 직접 다운로드해 확인했다.

\begin{longtable}{L{3.4cm}L{2.6cm}L{3.6cm}L{5.4cm}}
\caption{\texttt{geo\_id="KOR"} 실측 결과 (2026년 4--5월, \texttt{soc\_occupation} 카테고리)}\label{tab:kor}\\
\toprule
\textbf{hierarchy\_level} & \textbf{의미} & \textbf{확인된 행 수} & \textbf{제공 지표}\\
\midrule
\endfirsthead
\multicolumn{4}{l}{\small (표 \thetable\ 계속)}\\
\toprule
\textbf{hierarchy\_level} & \textbf{의미} & \textbf{확인된 행 수} & \textbf{제공 지표}\\
\midrule
\endhead
\bottomrule
\endfoot

0 (Detailed Occupation, SOC 6자리) &
세부직업 &
\texttt{pct} 813행, 시간 지표 \textbf{0행} &
\texttt{pct}만(문서 예측과 정확히 일치)\\

1 (Major Group, SOC 2자리) &
대분류(22개) &
전 지표 2,244행(22개 대분류 $\times$ 2개월 $\times$ 약 51개 지표) &
\texttt{human\_only\_time\_mean}$\cdot$\texttt{human\_with\_ai\_time\_mean} 등 전 지표\\

(참고) \texttt{overall} 카테고리 &
국가 전체 &
다수 &
전 지표(직업 미분해)\\
\end{longtable}

한국은 \texttt{geo\_id="KOR"}, \texttt{geo\_level="country"} 전체로 11,601행이 존재해 국가 포함 임계치를 통과한다. 대분류(레벨 1) 22개 전부가 관측되며, 모두 \texttt{human\_only\_time\_mean}(시간 단위)$\cdot$\texttt{human\_with\_ai\_time\_mean}(분 단위) 쌍을 갖는다.

\textbf{계산 예시}(2026년 5월, "Computer and Mathematical" 대분류, SOC 코드 15):

\[
\text{human\_only\_time\_mean} = 5.28\text{시간}, \quad \text{human\_with\_ai\_time\_mean} = 42.12\text{분}
\]
\[
\text{시간절감률} = 1 - \frac{42.12 / 60}{5.28} = 1 - 0.133 \approx 86.7\%
\]

이는 \texttt{08\_AEI데이터\_노동시간\_사용량\_스키마.tex} \S2.4가 "문서에 단위 미명시"로 남겨둔 문제를 해소하는 근거이기도 하다---V6 \texttt{data\_documentation.md}는 \texttt{human\_only\_time\_mean}의 단위를 \textbf{시간(hours)}, \texttt{human\_with\_ai\_time\_mean}의 단위를 \textbf{분(minutes)}으로 명시적으로 기술한다. \citet{tamkin-2025-estimating-ai-productivity-gains-from}의 원 방법론(전자는 시간 단위, 후자는 분 단위로 프롬프트 설계)이 공식 데이터셋에 그대로 반영되어 있음이 V6 문서로 확인된다.

% ============================================================
\section{GLOBAL 시간절감률 + KOR 사용량 결합: 기술적 검증과 방법론적 한계}
\label{sec:hybrid}
% ============================================================

\subsection{기술적 조인 가능성: 실측 결과}

\texttt{geo\_id="KOR"}의 세부직업(레벨 0) 사용 비중과 \texttt{geo\_id="GLOBAL"}의 세부직업 시간절감률을 O*NET-SOC 코드(\texttt{node\_external\_id})로 조인할 수 있는지를 실제로 확인했다.

\begin{longtable}{L{10.0cm}L{3.0cm}}
\caption{GLOBAL--KOR 조인 커버리지 (2026년 5월 기준)}\label{tab:coverage}\\
\toprule
\textbf{대상} & \textbf{코드 수}\\
\midrule
\endfirsthead
\multicolumn{2}{l}{\small (표 \thetable\ 계속)}\\
\toprule
\textbf{대상} & \textbf{코드 수}\\
\midrule
\endhead
\bottomrule
\endfoot

\texttt{geo\_id="KOR"}에서 관측된 세부직업 코드(레벨 0, \texttt{pct} 존재) & 424개\\
\texttt{geo\_id="GLOBAL"}에서 시간 지표가 존재하는 세부직업 코드(레벨 0) & 717개\\
한국 424개 코드 중 GLOBAL 시간 지표에 매칭되지 않는 코드 & \textbf{0개}\\
\end{longtable}

한국에서 관측되는 세부직업 코드 424개 전부가 GLOBAL 세부직업 시간 지표에 존재해, \textbf{커버리지 100\%로 조인 가능}하다. 즉 기술적으로는 다음 방식이 가능하다.

\[
\text{한국 직업 }o\text{의 시간절감률} := 1 - \frac{\text{GLOBAL\_human\_with\_ai\_time\_mean}_o / 60}{\text{GLOBAL\_human\_only\_time\_mean}_o}, \qquad \text{가중치} := \text{KOR\_pct}_o
\]

\subsection{방법론적 한계: 무엇을 가정하는 것인가}

이 결합이 기술적으로 가능하다는 것과, 결과가 "한국 고유의 실현 크기"를 나타낸다는 것은 다른 문제다. \texttt{08\_AEI데이터\_노동시간\_사용량\_스키마.tex} \S5.3이 정리한 "이론적 상한 $\times$ 실현 여부 $\times$ 실현 크기" 구조를 빌리면, 이 방식은 세 번째 항(실현 크기, 시간절감률)을 한국 고유값 대신 \textbf{GLOBAL 대리값(proxy)으로 치환}하는 것이다. 다음 세 가지를 방법론 절에 명시적 가정으로 적어야 한다.

\begin{enumerate}[label=(\arabic*)]
  \item \textbf{국가 간 동질성 가정}: 특정 직업의 AI 사용에 따른 시간절감 정도가 국가와 무관하게 동일하다고 가정한다. 이는 검증되지 않은 가정이며, 대분류(레벨 1) 수준에서는 한국 고유값이 실제로 존재하므로(\S\ref{sec:kor}), 이를 GLOBAL 대분류값과 비교하는 민감도 분석으로 가정의 타당성을 부분적으로 점검할 수 있다.
  \item \textbf{GLOBAL의 실질적 구성}: \texttt{GLOBAL}은 전세계 대화의 합산이므로, Claude 사용자 지리 분포상 영어권(특히 미국) 비중이 상당할 가능성이 크다. "GLOBAL을 한국에 적용한다"는 것의 실질적 의미는 "영어권 중심 표본의 직업별 시간절감 패턴을 한국에 적용한다"에 가까울 수 있다.
  \item \textbf{사용량과 시간절감률의 결합 층위 불일치}: 사용 비중(\texttt{pct})은 한국 고유 관측치이고 시간절감률은 전세계 관측치이므로, 산출되는 지표는 순수한 "관측된 실현 크기"가 아니라 "한국의 실현 여부 $\times$ 전세계 평균 실현 강도"라는 혼합 지표임을 논문 본문에서 구분해 표기해야 한다(예: "한국 적용 추정치"와 "한국 고유 관측치"를 별도 계열로 병기).
\end{enumerate}

% ============================================================
\section{SOC 코드는 몇 자리까지 가능한가}
\label{sec:link}
% ============================================================

\texttt{node\_external\_id}는 O*NET-SOC 2019 확장 코드(예: \texttt{11-3071.04}) 형식이다. 앞 6자리(\texttt{11-3071})가 표준 SOC \textbf{세부직업(Detailed Occupation) 코드}이며, 소수점 뒤 2자리는 SOC보다 더 세분화된 O*NET 고유 확장(하나의 SOC 코드가 여러 O*NET 직업으로 더 쪼개지는 경우에 부여됨)이다. \texttt{soc\_occupation} 카테고리는 이 중 두 계층만 노출한다---레벨 0(세부직업, 사실상 SOC 6자리)과 레벨 1(대분류, SOC 2자리)이며, 표준 SOC 체계의 중간 계층(3자리 소분류, 5자리 광의직업)은 AEI 데이터에 별도로 존재하지 않는다.

\begin{longtable}{L{3.8cm}L{2.6cm}L{4.6cm}L{4.6cm}}
\caption{SOC 계층별 국가(KOR) vs. GLOBAL 제공 지표}\label{tab:socdigit}\\
\toprule
\textbf{구분} & \textbf{SOC 자릿수} & \textbf{국가(KOR) 단위 제공 지표} & \textbf{GLOBAL 단위 제공 지표}\\
\midrule
\endfirsthead
\multicolumn{4}{l}{\small (표 \thetable\ 계속)}\\
\toprule
\textbf{구분} & \textbf{SOC 자릿수} & \textbf{국가(KOR) 단위 제공 지표} & \textbf{GLOBAL 단위 제공 지표}\\
\midrule
\endhead
\bottomrule
\endfoot

레벨 1: Major Group & 2자리 & 전 지표(시간 포함) & 전 지표(시간 포함)\\
레벨 0: Detailed Occupation & 6자리(O*NET 확장 시 6자리+2자리) & \texttt{pct}만 & 전 지표(시간 포함)\\
\end{longtable}

따라서 \S\ref{sec:hybrid}의 GLOBAL--KOR 결합 방식을 쓰면 \textbf{SOC 6자리(세부직업) 수준까지} 분석이 가능해진다. 이 방식을 쓰지 않고 한국 고유값만 쓴다면 \textbf{SOC 2자리(대분류, 22개 범주) 수준}으로 제한된다.

% ============================================================
\section{결론: 실행 권고}
\label{sec:conclusion}
% ============================================================

\begin{longtable}{L{4.6cm}L{10.6cm}}
\caption{질의별 결론 요약}\label{tab:summary}\\
\toprule
\textbf{질의} & \textbf{답}\\
\midrule
\endfirsthead
\multicolumn{2}{l}{\small (표 \thetable\ 계속)}\\
\toprule
\textbf{질의} & \textbf{답}\\
\midrule
\endhead
\bottomrule
\endfoot

어떤 파일을 쓸 것인가 & \texttt{release\_2026\_06\_26/data/aei\_claude\_ai\_2026-06-26.csv}(V6) 단일 파일. V4$\cdot$V5는 raw 스냅샷만 공개되어 있어 부적합\\

한국 고유값만으로 가능한 최대 해상도 & SOC 대분류(2자리, 22개), 직업별 시간절감률$\cdot$사용량 모두 가능\\

GLOBAL 시간절감률 + KOR 사용량 결합 시 최대 해상도 & SOC 세부직업(6자리, 한국에서 관측되는 424개 코드 전부 GLOBAL과 매칭 확인)\\

결합 방식의 조건 & 국가 간 시간절감 동질성 가정을 방법론 절에 명시. 대분류 수준에서 KOR 고유값과 GLOBAL값의 격차를 민감도 분석으로 함께 제시 권고\\
\end{longtable}

\texttt{code/01\_import\_aei\_kor.do}를 작성할 경우 다음을 주석으로 명시해야 한다(재현성 원칙, \texttt{CLAUDE.md}).

\begin{enumerate}[label=(\arabic*)]
  \item 원자료 출처: \texttt{release\_2026\_06\_26/data/aei\_claude\_ai\_2026-06-26.csv}(2026년 9월 시점 조회, 2026년 4--5월 데이터).
  \item 국가 필터 키: \texttt{geo\_id=="KOR"}(사용량), \texttt{geo\_id=="GLOBAL"}(세부직업 시간절감률 대리값 사용 시).
  \item 직업 코드 조인 키: \texttt{node\_external\_id}(O*NET-SOC 코드, 레벨 0) 또는 SOC 2자리(레벨 1, \texttt{node\_external\_id}가 대분류 코드).
  \item 시간 지표 단위: \texttt{human\_only\_time\_mean}=시간, \texttt{human\_with\_ai\_time\_mean}=분(V6 문서에 명시, \S\ref{sec:kor} 참조)---단위 환산 없이 결합하면 시간절감률이 오산출되므로 do 파일에서 반드시 분 단위로 통일 후 계산.
  \item 대분류(레벨 1) 한국 고유값과 GLOBAL 값의 비교표를 부속 산출물로 함께 저장해, \S\ref{sec:hybrid}의 동질성 가정을 사후 점검할 수 있도록 함.
\end{enumerate}

% ============================================================
\section{보론(2026-09-27 추가): V5 원자료 재검토와 V5$\cdot$V6의 직업 분류 활용 수준 비교}
\label{sec:v5v6}
% ============================================================

이 절은 2026년 9월 27일에 추가한 것이다. 앞 절들의 본문은 수정하지 않았다. 2절과 \S\ref{sec:conclusion} 표의 V5 관련 판단(``pre-enrichment 원시본이라 부적합'')은 \textbf{이 절의 검증 결과로 정정한다.} 두 서술이 충돌하면 이 절을 따른다.

\subsection{V5 원자료 직접 검증}

로컬에 저장한 V5 파일(\texttt{release\_2026\_03\_24/data/aei\_raw\_claude\_ai\_2026-02-05\_to\_2026-02-12.csv}, 2026년 2월 5--12일 1주)을 직접 집계했다.

\begin{enumerate}[label=(\arabic*)]
  \item \textbf{공개 임계치는 이미 적용되어 있다.} 한국(\texttt{geo\_id="KR"}) 과업 중 가장 작은 칸이 15건이고, 임계치에 못 미치는 과업은 \texttt{not\_classified}로 합산되어 있다(한국 대화의 24.9\%, GLOBAL 2.7\%). 파일명의 \texttt{raw}는 O*NET$\cdot$SOC 계층이나 인구 자료를 붙이는 enrichment 이전이라는 뜻으로 보이며, 프라이버시 처리 이전 자료라는 뜻은 아닌 것으로 판단된다.
  \item \textbf{\textcite{fan-2026-aggregate-gains-from-ai}가 쓴 R5 자료와 같은 파일이다.} 원문 Box 2의 수치와 정확히 일치한다. ``modify existing software'' 과업 41,586건, 분류된 대화 972,636건($=$ 100만 $-$ \texttt{not\_classified} 27,364), AI 없이 3.54시간 $\to$ AI 활용 시 17.8분.
  \item \textbf{스키마가 V6와 다르다.}
  \begin{itemize}
    \item 건수(\texttt{*\_count})가 있다. V6에는 점유율만 있다.
    \item \texttt{soc\_occupation} facet이 없다. 직업 단위로 가려면 과업$\to$직업 연계가 필요하다.
    \item 국가 코드가 2자리다(\texttt{KR}).
    \item 과업별 절감시간(\texttt{onet\_task::human\_only\_time} 등)은 GLOBAL에만 있다(3,259개 과업). 한국 행에는 국가 전체 평균만 있다.
  \end{itemize}
  \item \textbf{과업$\to$직업 연계가 사실상 완전하다.} \texttt{release\_2025\_09\_15/data/intermediate/onet\_task\_statements.csv}(O*NET 20.1, SOC 2010 코드)와 과업 텍스트를 대소문자 구분 없이 맞추면 다음과 같다.
  \begin{itemize}
    \item GLOBAL: 3,258개 과업 전부 매칭(\texttt{none}, \texttt{not\_classified} 제외)
    \item 한국: 299개 과업 전부 매칭
    \item 86개 과업(대화의 4.5\%)은 여러 직업에 걸려 있어, 해당 직업에 똑같이 나눈다.
    \item 연계 결과 세부직업 수는 GLOBAL 566개, 한국 171개다.
  \end{itemize}
\end{enumerate}

따라서 V5는 \textbf{분석에 쓸 수 있는 공개 집계본}이다. 건수와 과업별 95\% 신뢰구간이 있어, LCE 재현(Fan and Nguyen의 상한$\cdot$하한 시나리오 포함)에는 오히려 V6보다 유리하다.

\subsection{V6 한국 행의 추가 확인}

같은 기준으로 V6 한국(\texttt{geo\_id="KOR"}) 행을 다시 집계했다(2026년 4월 / 5월).

\begin{longtable}{L{9.0cm}L{2.8cm}L{2.8cm}}
\caption{V6 한국 행의 직업$\cdot$과업 단위 공개 현황}\label{tab:v6kor}\\
\toprule
항목 & 2026.4 & 2026.5 \\
\midrule
\endhead
\texttt{soc\_occupation} 레벨 0(세부직업) 공개 코드 수 & 389 & 424 \\
세부직업 \texttt{pct} 합계 & 96.16\% & 98.08\% \\
\texttt{onet} 레벨 0(과업) \texttt{pct} 합계 & 78.51\% & 81.32\% \\
세부직업 \texttt{pct} 최솟값 & 0.01\% & 0.01\% \\
국가 단위 직업별 절감시간 & 대분류 22개만 & 대분류 22개만 \\
\bottomrule
\end{longtable}

같은 한국 행인데도 과업 단위 합계(79--81\%)가 세부직업 단위 합계(96--98\%)보다 훨씬 낮다. 이는 V6가 공개 임계치를 \textbf{직업 단위 칸}에 따로 적용하기 때문으로 보인다. 직업은 여러 과업을 모은 단위라 임계치에 걸려 빠지는 몫이 적다. 반면 V5는 과업 단위로만 임계치가 적용되어 있어서, 과업을 직업으로 다시 묶어도 이미 빠진 몫(\texttt{not\_classified} 24.9\%)은 되살릴 수 없다.

V6에는 건수가 없고, 월별 표본 크기도 원자료 문서와 Cadences 보고서\parencite{massenkoff-2026-cadences} 어디에도 나오지 않는다. \texttt{pct}는 소수 둘째 자리로 반올림되어 있다. 따라서 V6 세부직업 값의 \textbf{포괄률은 높지만 정밀도(칸별 오차)는 평가할 수 없다.}

\subsection{V5$\cdot$V6 직업 분류 활용 수준 비교}

\begin{longtable}{L{3.4cm}L{5.6cm}L{5.6cm}}
\caption{용도별 V5$\cdot$V6 직업 분류 활용 수준}\label{tab:v5v6}\\
\toprule
용도 & V5 (2026.2, 1주) & V6 (2026.4--5, 월별) \\
\midrule
\endhead
\textbf{GLOBAL 직업별 절감시간} & 세부직업 가능 (과업 연계로 566개, 대화의 93\% 포괄, 건수$\cdot$신뢰구간 있음) & 세부직업 가능 (자체 제공 717개) \\
\textbf{KR 직업별 사용량} & \textbf{대분류 권장} (세부 171개로 나눌 수는 있으나 포괄률 70.3\%, 과업 299개 중 145개가 30건 미만) & \textbf{세부직업도 활용 가능} (389--424개, 포괄률 96--98\%). 단 정밀도는 확인할 수 없음 \\
\textbf{KR 직업별 절감시간} & 관측 불가 (국가 전체 평균만 있음) & \textbf{대분류(22개)만} 직접 관측 \\
\bottomrule
\end{longtable}

해석은 다음과 같다.
\begin{enumerate}[label=(\arabic*)]
  \item \textbf{GLOBAL에서는 두 릴리스 모두 세부직업까지 쓸 수 있다.} V5도 과업$\to$SOC 연계로 세부직업 수준을 충분히 만들 수 있고, 건수와 신뢰구간이 있어 더 투명하다.
  \item \textbf{한국 사용량은 V6가 세부 분류에 더 적합하다.} 이유는 표~\ref{tab:v6kor} 아래에 적은 임계치 적용 단위 차이다.
  \item \textbf{V6 세부직업 값의 정밀도는 보장되지 않는다.} 세부직업 결과는 순위나 상위 직업 위주로 해석한다. 회귀 등에 쓸 때는 4$\cdot$5월을 합치거나, SOC 소분류(3자리) 또는 ISCO 2자리로 묶은 버전을 함께 제시한다.
  \item \textbf{한국 직업별 절감시간은 두 릴리스 모두 세부직업 수준에서 관측되지 않는다.} 세부직업으로 가려면 어느 릴리스든 GLOBAL 과업$\cdot$직업별 절감시간을 빌려 와야 한다(\S\ref{sec:hybrid}의 결합 방식). 한국 고유 값은 V6 대분류뿐이다.
\end{enumerate}

요약하면 다음과 같다. \textbf{한국 사용량은 V5 대분류, V6 세부직업(해석 주의)까지 쓸 수 있다. 절감시간은 GLOBAL에서 두 릴리스 모두 세부직업까지, 한국에서는 V6 대분류까지만 쓸 수 있다.}

\subsection{참고: V5 SOC 대분류 잠정 집계}

아래는 scratchpad에서 awk로 계산한 \textbf{잠정치}다. do 파일로 재현하기 전이므로 인용하지 않는다. 절감시간은 과업별 (AI 없이 시간 $-$ AI 활용 분/60)을 대화 건수로 가중 평균한 값이고, 여러 직업에 걸린 과업은 똑같이 나눴다.

\begin{longtable}{L{1.0cm}L{5.4cm}L{2.6cm}L{3.0cm}L{2.6cm}}
\caption{V5 SOC 대분류별 대화 비중과 절감시간(잠정)}\label{tab:v5major}\\
\toprule
SOC & 직업군 & GLOBAL 대화 비중 & GLOBAL 대화당 절감시간 & KR 대화 비중 \\
\midrule
\endhead
15 & Computer \& Mathematical & 32.3\% & 3.58시간 & 35.0\% \\
25 & Educational Instruction \& Library & 13.5\% & 3.00 & 13.1\% \\
27 & Arts, Design, Entertainment, Sports \& Media & 10.4\% & 2.92 & 15.9\% \\
43 & Office \& Administrative Support & 9.4\% & 1.66 & 9.4\% \\
19 & Life, Physical \& Social Science & 6.3\% & 3.35 & 5.4\% \\
41 & Sales \& Related & 4.9\% & 1.32 & 4.3\% \\
\bottomrule
\end{longtable}

비중의 분모는 과업이 직업으로 매칭된 대화다(GLOBAL 929,714건, KR 20,429건).

\subsection{\S\ref{sec:conclusion} 실행 권고에 대한 보완}

\begin{itemize}
  \item \textbf{파일 선택}: V6 단일 파일이라는 권고를 ``V6를 기본으로 하되, V5(\texttt{release\_2026\_03\_24})도 사용 가능''으로 보완한다. 특히 건수가 필요한 분석(LCE 재현, 신뢰구간)에는 V5가 적합하다.
  \item \textbf{V5 do 파일 주석에 명시할 사항}: (1) 국가 코드 \texttt{KR}(2자리), (2) 과업$\to$직업 연계표 \texttt{release\_2025\_09\_15/data/intermediate/onet\_task\_statements.csv}(O*NET 20.1)와 매칭 키(소문자 변환한 과업 텍스트), (3) 다중직업 과업의 균등 배분 규칙, (4) 시간 단위(\texttt{onet\_task\_human\_only\_time\_mean}=시간, \texttt{onet\_task\_human\_with\_ai\_time\_mean}=분).
  \item \textbf{V5와 V6의 시계열 비교}: 두 릴리스는 기간 단위(1주 vs 월), 직업 코드 체계(연계로 만든 SOC 2010 vs V6 자체 제공 코드), 임계치 적용 단위가 다르다. 직업별 비중의 변화를 해석할 때 이 차이를 함께 적는다.
  \item \textbf{1P API}: V5 API 표본에는 Claude Code가 포함되어 있고\parencite[p.5]{massenkoff-2026-learning-curves}, V6 API에서는 제외되어 있다. V5와 V6의 API 직업 구성은 직접 비교하지 않는다.
\end{itemize}

% ============================================================
\printbibliography[title={참고문헌}]
% ============================================================

\end{document}
